
% ==============================================================================
% Fichier     : 02_fondamentaux.tex
% Description : Définitions, Concepts et Bases, Audit, Contrôle et SI
% Cours       : Audit des Systèmes d'Information (GES4062)
% Auteur      : MINKA MI NGUIDJOÏ Thierry Emmanuel
% Note        : Tous les concepts de la V0 sont repris intégralement.
% ==============================================================================

\chapter{Définitions, Concepts et Bases de l'Audit et du Contrôle}%
\label{chap:02_fondamentaux}

\begin{boxobjectifs}
\textbf{Objectif du chapitre:} Poser les fondations théoriques nécessaires à la
compréhension approfondie des systèmes d'information, de leur audit, et des enjeux
qui y sont associés. Une maîtrise rigoureuse de ces concepts est indispensable pour
appréhender les méthodologies, normes et pratiques qui guident l'évaluation, le
contrôle et la sécurisation des systèmes d'information.
\end{boxobjectifs}

% ==============================================================================
\section{Définitions, Concepts et Bases de l'Audit et du Contrôle}
% ==============================================================================

% ------------------------------------------------------------------------------
\subsection{Les Différences entre l'Audit et le Contrôle}
% ------------------------------------------------------------------------------

\begin{boxdef}{Audit et Contrôle}
\textbf{L'audit} a pour sujet un \textbf{processus}, tandis que le \textbf{contrôle}
s'applique à une \textbf{opération}.

Le but de l'audit est l'\textbf{optimisation du processus} sous revue pour en améliorer
la participation dans la \textbf{création de valeur} pour l'entité. Celui du contrôle
peut être de s'assurer de l'\textbf{exactitude}, la \textbf{complétude},
l'\textbf{existence}, la \textbf{cohérence}, \dots de l'opération visée.
\end{boxdef}

\begin{boxlizbiz}{Illustration Lizbiz's}
\begin{itemize}
  \item \textbf{Contrôle:} Le comptable vérifie que le montant d'une facture
    correspond à la commande reçue, il s'assure de la cohérence d'une opération.
  \item \textbf{Audit:} L'auditeur analyse l'ensemble du processus
    « Achat–Fournisseur » pour évaluer si les contrôles en place sont efficaces et
    systématiquement appliqués, et propose des optimisations créatrices de valeur.
\end{itemize}
\end{boxlizbiz}

% ------------------------------------------------------------------------------
\subsection{L'Assurance Raisonnable}
% ------------------------------------------------------------------------------

\begin{boxdef}{Assurance Raisonnable}
La notion d'assurance raisonnable est liée à la collecte des éléments probants
nécessaires pour que l'auditeur puisse conclure que les éléments analysés
\textbf{pris dans leur ensemble} ne comportent pas d'inexactitudes importantes,
pouvant entacher \textbf{les choix du commanditaire}.
\end{boxdef}

L'auditeur ne peut certifier à cent pour cent l'absence d'erreur ou de fraude, les
contraintes d'échantillonnage, de temps et la possibilité de dissimulation délibérée
constituent des limites inhérentes à l'exercice. L'assurance est donc raisonnable,
non absolue.

% ------------------------------------------------------------------------------
\subsection{Les Risques}
\label{subsec:Risques}
% ------------------------------------------------------------------------------

\subsubsection{Définition Académique du Risque}

\begin{boxdef}{Risque, Définition académique}
Le \termecle{risque} est la \textbf{probabilité} qu'une menace exploite une
vulnérabilité pour générer un \textbf{impact} sur une ressource identifiée.

L'élément capital pour l'existence du risque est donc la \textbf{ressource}.

Le risque est essentiellement probable; sa valeur est comprise dans l'intervalle
$]0,1[$. Dès matérialisation de la chose projetée, c'est-à-dire lorsque la valeur
vaut $1$, on parle d'\textbf{incident}. En revanche, lorsqu'elle vaut $0$, on parle
d'\textbf{opportunité}.
\end{boxdef}

\subsubsection{Généralités sur les Risques}

Les éléments minimaux constitutifs du risque sont :
\begin{itemize}
  \item la ressource ;
  \item la ou les vulnérabilités ;
  \item la ou les menaces ;
  \item le ou les impacts.
\end{itemize}

À ceux-ci peuvent s'ajouter : le ou les agents de menace, la compétence de l'agent
de menace, la motivation de l'agent de menace, etc.

L'évaluation des risques est un processus \textbf{cyclique} qui consiste à identifier
les risques, à préciser leur impact potentiel sur l'entreprise, à recenser les contrôles
en place pour les réduire, à évaluer la pertinence de ces contrôles et à déterminer si
les risques résiduels sont acceptables par l'entreprise.

\begin{boiteRemarque}{À retenir}
Par définition, sans ressource, pas de risque. Le gestionnaire de risque doit donc
savoir sur quels éléments il a la main, afin de bien dimensionner ses efforts.
\end{boiteRemarque}

\subsubsection{Traitement du Risque}

Il existe seulement \textbf{quatre méthodes} pour traiter le risque :

\begin{enumerate}
  \item \textbf{L'acceptation :} consiste à ne rien faire bien qu'étant pleinement
    conscient du risque encouru ;
  \item \textbf{L'évitement} (ou arrêt de l'activité) : préconise l'extinction de la
    raison du risque ;
  \item \textbf{Le transfert ou partage :} l'entité, par le mécanisme contractuel,
    s'assure de l'intervention d'un tiers pour l'aider le moment échéant à faire
    face à l'impact ;
  \item \textbf{La réduction ou mitigation :} se matérialise par le déploiement de
    contre-mesures ou l'amélioration de mécanismes de contrôle déjà en place pour
    agir soit sur la vulnérabilité, soit sur la menace, soit sur l'impact, ou une
    combinaison des trois.
\end{enumerate}

À l'issue de l'exécution d'une méthode de traitement, le risque résiduel doit avoir
été minimisé. Par abus de langage, quand on parle de \textbf{gestion des risques},
on entend généralement que la méthode de traitement est déjà choisie et qu'il s'agit
de la réduction ou mitigation.

\subsubsection{Principaux Groupes de Risques Informatiques}

Il y a trois principaux groupes de risques en informatique :
\begin{itemize}
  \item les risques \textbf{opérationnels} ;
  \item les risques \textbf{financiers} ;
  \item les risques \textbf{légaux} ou de non-conformité.
\end{itemize}

\begin{table}[htbp]
\centering
\caption{Quelques risques informatiques et leurs facteurs}
\label{tab:risques_informatiques}
\small
\begin{tabular}{p{6cm}p{7.5cm}}
\toprule
\textbf{Risques informatiques} & \textbf{Facteurs de risques} \\
\midrule
L'inadéquation du SI avec la stratégie de l'entité et les besoins des utilisateurs &
  Absence de schéma directeur ; manque d'implication de la direction ; absence de
  gouvernance informatique ; manque d'implication des utilisateurs métiers ; absence
  d'analyse de la valeur des SI mis en place. \\
\midrule
L'incapacité de l'organisation à redémarrer les systèmes en cas d'arrêt ou de destruction &
  Absence de sauvegarde régulière (sauvegardes externes) ; absence de plan de
  secours ; absence de site de secours. \\
\midrule
La sécurité du SI inadaptée au niveau de risque identifié et accepté par la direction &
  Absence de politique de sécurité ; failles en sécurité physique, logique, réseau,
  exploitation, postes et données. \\
\midrule
L'accès aux données et applications par des personnes non autorisées &
  Absence de gestion rigoureuse d'attribution des droits ; gestion insuffisante des
  mots de passe ; systèmes ne permettant pas une gestion fine des accès. \\
\midrule
Les applications informatiques non fiables &
  Erreurs de programmation par rapport aux spécifications ; applications
  insuffisamment testées ; utilisateurs insuffisamment impliqués dans le
  développement. \\
\midrule
L'indisponibilité du système informatique &
  Mauvaise gestion de l'environnement matériel (énergie, climatisation) ; absence
  de convention de service ; absence d'outil de surveillance ; absence de cellule
  réactive ; absence de contrat de maintenance. \\
\midrule
La mauvaise utilisation du système informatique &
  Applications non conviviales ; utilisateurs insuffisamment formés ;
  documentation insuffisante ; manque de contrôles bloquants. \\
\midrule
La non-conformité du système avec la législation &
  Politiques non conformes à la réglementation en cybersécurité, données
  personnelles, cybercriminalité, communications électroniques. \\
\midrule
La non-pérennité du système informatique &
  Utilisation de la sous-traitance sans transfert de compétence ; documentation
  inexistante ou non mise à jour ; obsolescence technologique ; forte dépendance
  vis-à-vis de personnes clés. \\
\bottomrule
\end{tabular}
\end{table}

\subsubsection{Risques en Audit}

L'informatique ne suffit pas à éliminer les risques, mais constitue un outil important
pour mieux les gérer et assurer la synergie de tous les secteurs de l'entreprise. Dans
le cadre des travaux d'audit, il faut obtenir une assurance raisonnable que les états
financiers sont exempts d'inexactitudes importantes.

Un audit ne nécessite pas la validation de la totalité des transactions ; l'exhaustivité
est difficilement atteignable, d'où l'importance de la gestion des risques en audit.

\paragraph{Risque Inhérent}
Le \termecle{risque inhérent} est le risque lié soit à l'activité, soit à l'entité,
soit au secteur dans lequel l'entité opère.

\paragraph{Risque de Non-contrôle}
Le \termecle{risque de non-contrôle} est fonction de l'efficacité avec laquelle la
conception et le fonctionnement du contrôle interne permettent d'atteindre les
objectifs de l'entité. Il subsiste toujours un risque de non-contrôle en raison des
limites liées au contrôle interne.

\paragraph{Risque de Non-détection}
Le \termecle{risque de non-détection} est le risque que l'auditeur ne détecte pas une
inexactitude qui pourrait être importante, soit isolément, soit cumulée avec d'autres
inexactitudes. Il est fonction de l'efficacité des procédés d'audit et de leur mise en
œuvre : il est donc lié à l'auditeur, à son expertise, son expérience, sa
compréhension du sujet d'audit et le choix de ses outils.

\paragraph{Risque Technologique}
Le \termecle{risque technologique} est lié à l'utilisation des technologies dans le
but d'atteindre les objectifs de l'entreprise. Il peut être pris en compte dans le
risque de non-détection.

\paragraph{Risque de Mission}
Le \termecle{risque de mission} est le risque que l'auditeur exprime une opinion
inappropriée sur les faits examinés contenant des inexactitudes significatives. Il est
représenté par la formule suivante :

\begin{boiteTheoreme}{Formule du Risque de Mission d'Audit}
\[
\mathbf{RMA} = RI \times RNC \times RND
\]
\begin{itemize}
  \item $\mathbf{RMA}$ : Risque de mission d'audit
  \item $\mathbf{RI}$ : Risque inhérent
  \item $\mathbf{RNC}$ : Risque de non-contrôle
  \item $\mathbf{RND}$ : Risque de non-détection
\end{itemize}
\end{boiteTheoreme}

\subsubsection{La Matrice des Risques}

\paragraph{Rôle et définition}
La \termecle{matrice des risques} est un outil opérationnel de gestion des risques.
Elle permet d'avoir sur le même support des informations utiles à la prise de décision
en rapport avec les risques. À ce titre, son usage est fortement recommandé par les
bonnes pratiques et les professionnels en la matière.

\paragraph{Contenu d'une Matrice des Risques}
La matrice des risques est un tableau récapitulatif qui met sur le même plan les points
suivants :
\begin{enumerate}
  \item un numéro d'ordre ;
  \item le domaine ou cycle concerné ;
  \item la structure compétente ;
  \item le poste de travail concerné ;
  \item le ou les titulaires du poste sur la période ;
  \item le logiciel utilisé ;
  \item la phase ou l'étape du processus visé ;
  \item le point de contrôle évalué ;
  \item le ou les critères d'évaluation ;
  \item la ou les questions d'évaluation ;
  \item la réponse associée ;
  \item le ou les éléments requis (futurs éléments probants) ;
  \item le risque associé ;
  \item la conséquence majeure en cas de réalisation du risque relevé ;
  \item la criticité subséquente ;
  \item un commentaire le cas échéant.
\end{enumerate}

\begin{boxexo}{Exercice}
Chaque apprenant proposera une matrice des risques dans un fichier Excel contenant
les seize éléments énumérés ci-dessus.
\end{boxexo}

% ------------------------------------------------------------------------------
\subsection{Le Contrôle Interne}
% ------------------------------------------------------------------------------

\subsubsection{Définition du Contrôle Interne}

\begin{boxdef}{Contrôle Interne, Définition COSO}
Le référentiel \norme{COSO} définit le contrôle interne comme un processus mis en
œuvre par les dirigeants qui couvre tous les niveaux de l'entreprise et est destiné à
fournir une \textbf{assurance raisonnable} quant à la réalisation des trois objectifs
suivants :
\begin{itemize}
  \item l'efficacité et l'efficience des opérations ;
  \item la fiabilité des informations financières ;
  \item la conformité aux lois et règlements.
\end{itemize}
\end{boxdef}

\subsubsection{Généralités sur le Contrôle Interne}

Le contrôle interne est une notion indispensable pour toutes les entreprises,
notamment au niveau managérial. Il s'agit d'un véritable système de sécurité/veille,
composé d'un ensemble de méthodes, de règles et de procédures définies par
l'entreprise, mais aussi d'un ensemble de postes de travail et de fonctions présentes
dans l'organigramme. Il s'adapte selon la taille et les caractéristiques de chaque
entité : plus la taille de l'entreprise est importante, plus les procédés du contrôle
interne devraient être conséquents.

Le contrôle interne vise à avoir la pleine maîtrise de son entreprise et à éviter les
erreurs. Il garantit la qualité de la gestion, du système mis en place et la bonne
utilisation des ressources. Pour le mettre en place, il est nécessaire de posséder une
vision claire de l'entité, ce qui passe par la bonne codification des processus de
production.

\subsubsection{Évaluation du Système de Contrôle Interne}

L'évaluation du système de contrôle interne a pour but de renseigner sur l'état du
contrôle interne mis en place par l'entité en lien avec le sujet d'audit, afin d'en
accroître la plus-value. Dans le cadre de l'audit d'un pan de l'entité, l'auditeur
devra réaliser l'abstraction du système de contrôle interne couvert par le mandat.

Le référentiel faisant autorité en matière d'évaluation de système de contrôle interne
est le \norme{COSO}, avec ses \textbf{cinq composantes} et ses
\textbf{dix-sept principes}, qui assure une couverture de bout en bout de l'entreprise.

\begin{boxexo}{Exercice}
Chaque apprenant proposera dans la première colonne les 5 composantes du COSO,
et dans la seconde les 17 principes rangés par composante. Une fois l'exercice
réalisé : dans quelle cellule du tableau ainsi constitué devrait être rangé les
bordereaux de transmission utilisés dans l'entité ?
\end{boxexo}

% ------------------------------------------------------------------------------
\subsection{La Règle des Quatre Yeux}
% ------------------------------------------------------------------------------

\begin{boiteTheoreme}{Règle des Quatre Yeux, Énoncé}
La première paire d'yeux observe ; la seconde peut \textbf{infirmer, confirmer ou
réformer}.
\end{boiteTheoreme}

Un auditeur n'achève jamais un travail seul, quel que soit son niveau d'expertise ou
d'expérience. Il a toujours besoin que son travail soit révisé : pour confirmer qu'il ne
s'est pas trompé, ou pour indiquer qu'il s'est trompé, ou encore pour mettre son
travail en perspective.

\begin{boxexo}{Exercice}
Chaque apprenant proposera deux exemples d'application de la règle des quatre yeux.
\end{boxexo}

% ------------------------------------------------------------------------------
\subsection{L'Allégorie des Deux Mains}
% ------------------------------------------------------------------------------

On peut se servir de ses deux mains pour illustrer ce qu'est l'audit.

\begin{itemize}
  \item Dans la \textbf{première main}, l'auditeur charge le
    \textbf{corpus réglementaire} dont il se servira.
  \item Dans la \textbf{seconde}, il place les \textbf{faits observés}, la pratique
    sur le terrain.
  \item Puis il place ses deux mains en regard pour en dégager les
    \textbf{écarts et différences} éventuels.
\end{itemize}

% ------------------------------------------------------------------------------
\subsection{Les Trois Piliers de l'Audit Interne}
% ------------------------------------------------------------------------------

L'audit interne repose sur trois piliers fondamentaux :

\begin{enumerate}
  \item \textbf{La spécification des tâches :} veut que ce qu'il y a à faire soit
    bien défini, précisé et communiqué. Elle se matérialise au niveau de l'agent par
    la \textbf{fiche de poste}.

  \item \textbf{La séparation des tâches et fonctions incompatibles :} a pour but
    d'éviter qu'un maillon de la chaîne ne concentre suffisamment de prérogatives
    pour, seul, pouvoir exécuter de bout en bout une procédure. Ce pilier est
    implémenté en finances publiques au Cameroun par le principe qui stipule qu'un
    seul individu ne peut pas cumuler à la fois les fonctions d'\textbf{ordonnateur}
    et de \textbf{comptable public} pour la même opération.

  \item \textbf{Le contrôle réciproque des tâches :} précise que le palier supérieur
    prolonge le travail du palier inférieur en le validant. Ce pilier se matérialise
    dans l'administration publique notamment par le système de \textbf{visa
    hiérarchique} sur les documents. C'est aussi le cas dans les divers \textit{workflows}
    où, sans la validation explicite du N+1, la procédure n'avance pas.
\end{enumerate}

% ------------------------------------------------------------------------------
\subsection{Point de Contrôle}
% ------------------------------------------------------------------------------

\begin{boxdef}{Point de Contrôle}
Un \termecle{point de contrôle} représente un élément dans la procédure auditée pour
lequel l'auditeur souhaiterait obtenir des éclaircissements, en questionnant la
conformité au référentiel établi et en récoltant les éléments d'analyse.

Dans la pratique, si l'on linéarise la procédure à auditer, chaque acteur ou changement
de phase peut être matérialisé par un point de contrôle.
\end{boxdef}

\begin{boxlizbiz}{Exemple, Procédure école}
Soit la procédure suivante : l'apprenant part du portail de l'école pour la salle de
classe. Les points de contrôle peuvent être :
\begin{itemize}
  \item le portail ;
  \item l'entrée dans le bâtiment ;
  \item l'entrée dans la salle de classe.
\end{itemize}
\end{boxlizbiz}

\begin{boxexo}{Exercice}
Chaque apprenant proposera un exemple de procédure et en identifiera les points de
contrôle. La procédure doit donner lieu à au moins quatre points de contrôle.
\end{boxexo}

% ------------------------------------------------------------------------------
\subsection{Objectif de Contrôle}
% ------------------------------------------------------------------------------

\begin{boxdef}{Objectif de Contrôle}
On peut définir un \termecle{objectif de contrôle} comme étant la raison pour laquelle
l'auditeur évalue un point de contrôle. Pour chaque point de contrôle, il y a au
moins un objectif de contrôle.
\end{boxdef}

\paragraph{Formulation}
Il est usuel de formuler les objectifs de contrôle de l'une des façons suivantes :
\begin{itemize}
  \item \emph{S'assurer que xxxxxxxxxxxxx ;}
  \item \emph{Vérifier que xxxxxxxxxxxxx.}
\end{itemize}

\begin{boxlizbiz}{Exemples, Procédure école}
Pour les points de contrôle identifiés précédemment :
\begin{itemize}
  \item \textbf{Le portail :} S'assurer que l'apprenant soit identifié au niveau du portail.
  \item \textbf{L'entrée dans le bâtiment :} S'assurer que seuls les apprenants en règle
    avec la scolarité accèdent au bâtiment.
  \item \textbf{L'entrée dans la salle de classe :} Vérifier qu'aucun apprenant n'entre en
    salle de classe après l'enseignant.
\end{itemize}
\end{boxlizbiz}

% ------------------------------------------------------------------------------
\subsection{Critère d'Évaluation}
% ------------------------------------------------------------------------------

\begin{boxdef}{Critère d'Évaluation}
Les \termecle{critères d'évaluation} sont la matérialisation des jalons dont
l'assemblage permet de réaliser l'objectif de contrôle. Il faut donc satisfaire
l'ensemble des critères d'évaluation d'un point de contrôle pour que son évaluation
soit concluante. Dans la pratique, il faut généralement plusieurs critères d'évaluation
par objectif de contrôle.
\end{boxdef}

\begin{boxlizbiz}{Exemple, Portail de l'école}
Pour l'objectif « S'assurer que l'apprenant soit identifié au niveau du portail »,
les critères d'évaluation sont :
\begin{itemize}
  \item L'apprenant a une pièce d'identité ;
  \item Un vigile est présent à la guérite du portail ;
  \item Le vigile dispose de la liste d'apprenants s'étant acquittés de leurs droits
    académiques ;
  \item La liste dont dispose le vigile est à jour ;
  \item La comparaison entre la pièce d'identité de l'apprenant et la liste à jour dont
    dispose le vigile est probante.
\end{itemize}
\end{boxlizbiz}

% ------------------------------------------------------------------------------
\subsection{Question d'Évaluation}
% ------------------------------------------------------------------------------

\begin{boxdef}{Question d'Évaluation}
Les \termecle{questions d'évaluation} permettent de mesurer l'atteinte des objectifs
de contrôle en transformant quasi-systématiquement chaque critère d'évaluation en
question libellée de façon univoque et de manière que les \textbf{réponses soient
booléennes}.
\end{boxdef}

\begin{boxlizbiz}{Exemple, Portail de l'école}
Pour l'objectif << S'assurer que l'apprenant soit identifié au niveau du portail >> :
\begin{itemize}
  \item L'apprenant a-t-il une pièce d'identité ?
  \item Y a-t-il un vigile à la guérite du portail ?
  \item Le vigile dispose-t-il d'une liste d'apprenants s'étant acquittés de leurs droits
    académiques ?
  \item Cette liste est-elle à jour ?
  \item Le vigile compare-t-il la pièce d'identité de l'apprenant à la liste à jour dont
    il dispose ?
\end{itemize}
\end{boxlizbiz}

% ------------------------------------------------------------------------------
\subsection{Constatation d'Audit}
% ------------------------------------------------------------------------------

\subsubsection{Définitions}

Plusieurs définitions coexistent dans la littérature ; au final elles renvoient toutes
à la même réalité :

\begin{boxdef}{Constatation d'Audit}
Une \termecle{constatation d'audit} est la consignation du fait observé par l'auditeur,
c'est-à-dire un \textbf{écart par rapport au référentiel} utilisé comme
\textit{baseline}.

Quoique peu courant, une constatation d'audit n'est pas toujours un écart
\textbf{négatif} : certaines entités performent au-delà de la baseline de référence —
il s'agit alors d'un écart \textbf{positif}.
\end{boxdef}

La définition consistant à définir la constatation d'audit comme une violation d'une
norme établie et diffusée n'a pas été retenue ici, pour cette raison.

\subsubsection{Exemples}

\paragraph{Constatation Positive}
L'équipe de mission a constaté qu'en sus de la règle établie et des bonnes pratiques
en la matière, qui fixent le nombre de missions d'audit annuelle pour les applications
en service à une (1), l'entreprise Lizbiz's en a réalisé trois (3) en 2022.

\paragraph{Constatation Négative}
L'équipe de mission a constaté qu'en violation du second pilier de l'audit interne qui
grave la séparation des tâches et fonctions incompatibles dans le marbre, repris et
surligné par le point 10.3.2 d'\norme{ISO 27002} qui interdit le cumul des fonctions
d'administrateur de base de données et celui d'administrateur réseau, dans l'entreprise
Lizbiz's, Monsieur \textsc{Bomba} a régulièrement exercé ces deux fonctions pendant
toute la période sous revue.

% ------------------------------------------------------------------------------
\subsection{Le Contradictoire}
% ------------------------------------------------------------------------------

\begin{boxdef}{Contradictoire}
Élément << sine qua non >> du droit de la défense, le \termecle{contradictoire}
est repris en audit pour donner la possibilité à l'audité de présenter sa version du
fait querellé, appuyée d'éléments probants suffisants au regard de la constatation
dressée par l'auditeur.

Il matérialise l'\textbf{obligation d'impartialité} et la \textbf{preuve d'humilité}
de l'auditeur, qui soumet à l'audité le premier jet de ce qu'il pense avoir établi comme
constat, avant sa consignation définitive dans le rapport.
\end{boxdef}

% ------------------------------------------------------------------------------
\subsection{Les Observations}
% ------------------------------------------------------------------------------

Dans certains types de rapport de mission, notamment ceux d'ISC comme le
\sigle{CONSUPE} (Contrôle Supérieur de l'État), l'observation est la forme arrêtée
pour consigner les travaux d'audit.

\begin{boxdef}{Structure d'une Observation}
Une observation se compose de :
\begin{enumerate}
  \item un numéro d'ordre ;
  \item l'énoncé de la norme visée ;
  \item la constatation d'audit ;
  \item les risques ou conséquences du fait observé ;
  \item le contradictoire ;
  \item l'analyse de l'auditeur ;
  \item la conclusion sur le fait observé ;
  \item la ou les recommandations pour contenir, réparer ou éviter à l'avenir la
    situation observée.
\end{enumerate}
\end{boxdef}

\begin{boiteRemarque}{Rappel}
Il faudra le cas échéant se conformer à la structure arrêtée par son entité, qui peut
différer de ce modèle générique.
\end{boiteRemarque}

% ------------------------------------------------------------------------------
\subsection{Sur Place et Sur Pièces}
% ------------------------------------------------------------------------------

\begin{boxdef}{Sur Place et Sur Pièces}
Ce concept renvoie au fait que, pour réaliser un audit, l'auditeur doit se déplacer
physiquement de manière à être présent sur les lieux, y recueillir les pièces
justificatives nécessaires, puis les traiter en personne et se prononcer sur leur
caractère suffisant ou non.

\textbf{L'audit ne se base donc pas sur des ouï-dire} ou des allégations non factuelles,
même si ces éléments peuvent être utilisés pour ouvrir ou fermer des pistes d'audit.
\end{boxdef}

% ------------------------------------------------------------------------------
\subsection{L'Aveu de Carence}
% ------------------------------------------------------------------------------

\begin{boxdef}{Aveu de Carence}
Lorsque l'audité est sollicité sur un fait pour en apporter une explication
complémentaire à celle de l'équipe de mission et ne se manifeste pas, cette dernière
peut conclure à l'\termecle{aveu de carence} après avoir fait la preuve que l'audité
a été relancé sur le fait querellé plusieurs fois, par tout moyen laissant trace, et n'a
pas réagi.

L'aveu de carence signifie l'\textbf{acceptation tacite} par l'audité du fait qui est
porté à son attention. C'est donc une \textbf{exception au principe du contradictoire}.
\end{boxdef}

% ==============================================================================
\section{Définitions, Concepts et Bases des Systèmes d'Information}
% ==============================================================================

% ------------------------------------------------------------------------------
\subsection{Notion de Système}
% ------------------------------------------------------------------------------

Plusieurs définitions coexistent dans la littérature, certaines à la limite de
l'antinomisme. Dans le cadre de ce cours, un \termecle{système} sera considéré comme
étant un \textbf{groupe d'éléments interagissant ensemble pour rendre un résultat
prédictible}.

% ------------------------------------------------------------------------------
\subsection{Les Systèmes d'Information}
% ------------------------------------------------------------------------------

\subsubsection{Système Informatique}

\begin{boxdef}{Système Informatique}
Représente l'ensemble des \textbf{logiciels et matériels informatiques} participant à
l'acquisition, au stockage, au traitement, au transport et à la diffusion de
l'information sous forme électronique au sein de l'organisation.
\end{boxdef}

\subsubsection{Système d'Information}

\begin{boxdef}{Système d'Information}
Ensemble des \textbf{ressources humaines, matérielles, logicielles et procédurales}
participant à l'acquisition, au stockage, au traitement, au transport et à la diffusion
de l'information sous toutes ses formes au sein de l'organisation.

Il s'agit d'un \textbf{système sociotechnique} composé de deux sous-systèmes :
l'un social (structure organisationnelle et personnes) et l'autre technique
(\textit{hardware}, \textit{software}, équipements de télécommunication et processus
d'affaires).
\end{boxdef}

\subsubsection{Fonction Informatique}

\begin{boxdef}{Fonction Informatique}
La \termecle{fonction informatique} comprend, outre le système informatique, les
personnes, processus, ressources financières et informationnelles. Elle a pour but de
fournir à l'entreprise les éléments nécessaires à la réalisation optimale de ses
missions.
\end{boxdef}

\subsubsection{Fonction d'Information}

\begin{boxdef}{Fonction d'Information}
Ensemble organisé de personnes, de procédures et d'équipement qui a pour objet de
réunir, de trier, d'analyser, d'évaluer et de distribuer, en temps utile, de l'information
pertinente et valide, provenant de sources internes et externes à l'organisation, afin
que tous ceux qui ont à prendre des décisions disposent des éléments leur permettant
de choisir l'action la plus appropriée au moment adéquat.
\end{boxdef}

\subsubsection{Problèmes Inhérents aux SI}

Comme tous les systèmes, les SI peuvent être sujets à des problèmes liés soit à
l'agencement des éléments qui les constituent, soit aux fonctions qu'ils fournissent,
ou une combinaison des deux. Tous les problèmes inhérents au SI peuvent donc en
absolu se classer dans l'une de ces deux catégories.

\subsubsection{Informatisation des SI}

Tous les systèmes d'information ne sont ni informatisés ni informatisables.
L'informatisation d'un système d'information est une phase critique de sa vie. Des
préalables doivent être respectés :
\begin{itemize}
  \item la définition claire et précise du but et des objectifs de l'informatisation ;
  \item les termes de référence de son informatisation, qui reprennent la définition du
    besoin, le périmètre des travaux, une certaine vision du futur système, les délais
    de réalisation, et les autres contraintes connexes ;
  \item l'allocation des ressources pour le projet d'informatisation, notamment
    matérielles, humaines et financières.
\end{itemize}

L'informatisation d'un système d'information est donc un projet à traiter comme tel,
dans le strict respect des usages et bonnes pratiques en la matière.

\subsubsection{Types de Migration/Changement de SI}

Il existe trois grands types de changement de système :

\begin{table}[htbp]
\centering
\caption{Différents types de migration de système}
\label{tab:migration}
\small
\begin{tabular}{clp{5cm}p{5cm}}
\toprule
\textbf{N°} & \textbf{Type} & \textbf{Avantages majeurs} & \textbf{Inconvénients majeurs} \\
\midrule
1 & Direct & Rapidité d'exécution. & Aucun retour arrière possible. \\
2 & Par étape &
  Possibilité de retour arrière si une étape se passe mal ; assimilation progressive
  du nouveau système brique par brique. &
  Temps relativement long de mise en œuvre ; ressource en infrastructure requise pour
  supporter les deux systèmes. \\
3 & En parallèle &
  Temps d'apprentissage suffisant du nouveau système ; comparaison poussée des
  deux systèmes ; possibilité de basculer vers le système le plus avantageux. &
  Ressources requises pour faire tourner les deux systèmes simultanément. \\
\bottomrule
\end{tabular}
\end{table}

Pour chacun de ces types de migration, l'auditeur devra recueillir la documentation et
observer comment et si les règles inhérentes à la migration de système ont été
respectées, notamment :
\begin{itemize}
  \item l'exhaustivité du transfert des données de l'ancien vers le nouveau système ;
  \item la convertibilité des données le cas échéant ;
  \item l'intégrité des données au moment de leur transfert ;
  \item la préservation de la confidentialité des données ;
  \item l'exactitude ou la convergence des résultats du nouveau système ;
  \item la préservation des principes de séparation des tâches et fonctions
    incompatibles.
\end{itemize}

% ------------------------------------------------------------------------------
\subsection{Système d'Exploitation}
% ------------------------------------------------------------------------------

\begin{boxdef}{Système d'Exploitation}
Un \termecle{système d'exploitation} est un logiciel constitué de plusieurs fonctions
dont le rôle est la gestion du matériel informatique qui constitue son hôte et les
logiciels qui y sont installés. Il est encore appelé \textbf{logiciel de base}, car les
autres logiciels sont installés par-dessus lui.
\end{boxdef}

Microsoft est le constructeur du célèbre système d'exploitation Windows, dont quelques
déclinaisons sont : Windows XP, Windows 10, Windows 11.

% ------------------------------------------------------------------------------
\subsection{Système de Fichier}
% ------------------------------------------------------------------------------

\begin{boxdef}{Système de Fichier}
Un \termecle{système de fichier} décrit la façon dont le système d'exploitation
organise la gestion des éléments sur un support physique (disque dur, disque
amovible, clé USB).
\end{boxdef}

\norme{FAT}, \norme{FAT32} et \norme{NTFS} sont des systèmes de fichiers Windows.

% ------------------------------------------------------------------------------
\subsection{Fichiers}
% ------------------------------------------------------------------------------

\begin{boxdef}{Fichier}
Un \termecle{fichier} est la matérialisation du stockage sur support physique des
données suivant un format défini dans le système de fichier. À chaque fichier est
associé un logiciel particulier pour sa lecture et modification ; le lien entre les deux
se fait par une \textbf{extension}.
\end{boxdef}

On peut classer les fichiers par leur taille, leur extension, leur nom ou leur date de
modification. À titre d'exemple, Word ouvre en standard les fichiers d'extension
\texttt{.doc} et \texttt{.docx} ; pour Excel, ce sont les extensions \texttt{.xls} et
\texttt{.xlsx}.

% ------------------------------------------------------------------------------
\subsection{Base de Données}
% ------------------------------------------------------------------------------

\begin{boxdef}{Base de Données}
Une \termecle{base de données} est une collection organisée de données structurées,
généralement stockées électroniquement dans un système informatique. On peut la
concevoir comme un classeur électronique qui contient des tiroirs et leurs contenus.
Les bases de données sont de plusieurs types, les plus répandus étant objet et
relationnel. C'est le \sigle{SGBD} (Système de Gestion de Base de Données) qui gère
et contrôle la base de données.
\end{boxdef}

% ------------------------------------------------------------------------------
\subsection{Réseau Informatique}
% ------------------------------------------------------------------------------

\begin{boxdef}{Réseau Informatique}
Un \termecle{réseau informatique} est un système constitué de matériel et logiciel
informatiques interagissant entre eux pour acheminer des données.
\end{boxdef}

% ------------------------------------------------------------------------------
\subsection{Sécurité Informatique}
% ------------------------------------------------------------------------------

\begin{boxdef}{Sécurité Informatique}
La \termecle{sécurité informatique} peut se définir comme un ensemble constitué de
moyens humains, matériels, procéduraux et financiers mis en œuvre dans le but de
garantir l'\textbf{intégrité}, la \textbf{disponibilité} et la \textbf{confidentialité}
des informations traitées, stockées et acheminées au moyen d'outils et de procédés
informatiques.
\end{boxdef}

% ==============================================================================
\section{Les Contrôles en Audit des Systèmes d'Information}
% ==============================================================================

% ------------------------------------------------------------------------------
\subsection{Rôle}
% ------------------------------------------------------------------------------

Simplement dit, le rôle d'un contrôle est de s'assurer du \textbf{respect d'une règle
connue, établie et précise}, avant la réalisation ou non d'une action subséquente.

% ------------------------------------------------------------------------------
\subsection{Les Contrôles dans les SI}
% ------------------------------------------------------------------------------

Lors de l'audit des SI, deux catégories majeures de contrôles sont évaluées :
les \textbf{contrôles généraux des systèmes d'information} et les
\textbf{contrôles d'application}. Ces deux types de contrôles agissent de manière
complémentaire mais divergent dans leur portée et leur approche.

% ------------------------------------------------------------------------------
\subsubsection{Contrôles Généraux}
% ------------------------------------------------------------------------------

\paragraph{Définition et objectif}

\begin{boxdef}{Contrôles Généraux (\textit{GITC})}
Les \termecle{contrôles généraux des systèmes d'information}
(\textit{General IT Controls}, \sigle{GITC}) concernent l'\textbf{environnement
informatique global} d'une organisation. Ils sont mis en place pour garantir que
l'infrastructure informatique est gérée de manière sécurisée et fiable. Ces contrôles
sont \textbf{transversaux} et s'appliquent à l'ensemble des systèmes et applications
informatiques, assurant ainsi une base stable pour leur bon fonctionnement.
\end{boxdef}

\begin{boxlizbiz}{Exemple Lizbiz's}
La société Lizbiz's utilise plusieurs serveurs pour héberger ses systèmes de gestion
agricole et commerciale. Les contrôles généraux vont s'assurer que ces serveurs sont
protégés contre les accès non autorisés, que les mises à jour de sécurité sont
régulièrement effectuées, et que des procédures de sauvegarde sont en place en cas de
défaillance matérielle ou d'incident informatique.
\end{boxlizbiz}

\paragraph{Quelques Champs d'Application}

\subparagraph{Sécurité des Systèmes et Gestion des Accès}
Ce contrôle vise à s'assurer que seuls les utilisateurs autorisés ont accès aux
systèmes d'information. Il implique des mesures de protection comme les mots de
passe robustes, l'authentification multi-facteurs (\sigle{MFA}), et une gestion stricte
des autorisations d'accès.

\begin{boxlizbiz}{Exemple Lizbiz's}
Lizbiz's met en place une politique de mots de passe obligeant les employés à changer
leur mot de passe tous les 90 jours, avec des exigences de complexité. L'accès aux
systèmes critiques nécessite une authentification par application mobile en plus du
mot de passe.
\end{boxlizbiz}

\subparagraph{La Gestion des Changements}
Les processus de gestion des changements garantissent que toutes les modifications
apportées aux systèmes informatiques (mises à jour logicielles, changements
d'infrastructure) sont correctement planifiées, testées et approuvées, minimisant les
risques de perturbation.

\subparagraph{La Gestion des Opérations Informatiques}
Les opérations quotidiennes des systèmes informatiques doivent être surveillées pour
assurer leur continuité, sauvegardes régulières, gestion des incidents, planification
de la maintenance.

\subparagraph{Plans de Continuité des Activités et Reprise après Sinistre}
Le premier prévoit un ensemble de mécanismes pour s'assurer que les services
identifiés comme critiques continuent d'être fonctionnels quel que soit l'incident.
Le second prévoit la restauration des systèmes d'information à l'état normal en cas
d'incident majeur. L'objectif commun : \textbf{minimiser les interruptions des activités
critiques}.

% ------------------------------------------------------------------------------
\subsubsection{Contrôles d'Application}
% ------------------------------------------------------------------------------

\paragraph{Définition et objectif}

\begin{boxdef}{Contrôles d'Application}
Les \termecle{contrôles d'application} sont une transcription informatique des
contrôles organisationnels. Encore appelés \textbf{contrôles programmés}, ils sont
encodés au moyen de programmes informatiques dans les logiciels pour matérialiser
l'existence d'une règle. Ils sont fréquents dans les opérations d'entrée, de traitement
et de sortie de l'application (\textbf{modèle IPO}). Leur objectif est de garantir
l'intégralité, la fiabilité et l'exactitude du traitement des données.
\end{boxdef}

\paragraph{Quelques Champs d'Application}

\subparagraph{Contrôles d'Entrée de Données}
Ces contrôles s'assurent que les données saisies dans les systèmes d'information
sont exactes, complètes et conformes aux règles et formats établis.

\begin{boxlizbiz}{Exemple Lizbiz's}
Dans le système de gestion des récoltes, un contrôle empêche l'enregistrement d'une
nouvelle livraison de coton si les informations sur le poids et la qualité ne sont pas
correctement renseignées.
\end{boxlizbiz}

\subparagraph{Contrôle de Traitement de Données}
Ces contrôles vérifient que les transactions sont correctement exécutées selon les
règles métier prédéfinies.

\begin{boxlizbiz}{Exemple Lizbiz's}
Dans le système de facturation, un contrôle empêche la génération d'une facture tant
que la livraison n'a pas été confirmée dans le système et que toutes les étapes de
validation n'ont pas été suivies.
\end{boxlizbiz}

\subparagraph{Contrôle de Sortie de Données}
Ces contrôles garantissent que les informations générées par les systèmes (rapports
financiers, états de gestion) sont exactes, complètes et conformes aux attentes.

\subparagraph{Contrôle d'Accès Spécifiques aux Applications}
Ces contrôles s'assurent que seules les personnes habilitées peuvent effectuer
certaines actions sensibles dans les applications.

\begin{boxlizbiz}{Exemple Lizbiz's}
Dans le système de gestion des ressources humaines, seuls les gestionnaires des paies
peuvent approuver les versements salariaux, chaque approbation étant enregistrée avec
un horodatage et l'identifiant de l'utilisateur.
\end{boxlizbiz}

% ------------------------------------------------------------------------------
\subsubsection{Complémentarité entre Contrôles Généraux et Contrôles d'Application}
% ------------------------------------------------------------------------------

\paragraph{Interdépendance fonctionnelle}
Les contrôles généraux créent un environnement sécurisé et stable qui permet aux
contrôles d'application de fonctionner efficacement. Un contrôle général de gestion
des accès garantit que seuls les utilisateurs autorisés peuvent interagir avec une
application spécifique ; un contrôle général de gestion des changements prévient les
erreurs systémiques susceptibles d'affecter les applications.

\paragraph{Sécurisation renforcée}
Les contrôles généraux protègent l'infrastructure et l'environnement global, tandis
que les contrôles d'application s'assurent que les processus métier spécifiques sont
correctement exécutés.

\begin{boxlizbiz}{Exemple, Vulnérabilité de complémentarité}
Lors d'un audit de Lizbiz's, il a été constaté que bien que le système de gestion des
ressources humaines ait des contrôles d'application solides pour valider les congés des
employés, une faiblesse dans les contrôles généraux de gestion des accès a permis à
un utilisateur non autorisé d'accéder au système et de modifier des données sensibles.
Cela démontre l'importance de la complémentarité des deux types de contrôles.
\end{boxlizbiz}

\paragraph{Prévention des erreurs à différents niveaux}
Les contrôles généraux sont préventifs à un niveau global ; les contrôles d'application
se concentrent sur la prévention des erreurs spécifiques aux transactions ou aux
processus métier.

% ------------------------------------------------------------------------------
\subsubsection{Divergences entre Contrôles Généraux et Contrôles d'Application}
% ------------------------------------------------------------------------------

\paragraph{Portée}
Les contrôles généraux s'appliquent à \textbf{l'ensemble des systèmes et des
infrastructures} informatiques de l'organisation. Les contrôles d'application, en
revanche, sont \textbf{spécifiques à des logiciels ou à des systèmes particuliers},
concentrés sur des processus métier bien définis.

\paragraph{Focalisation}
Les contrôles généraux se focalisent sur la gestion et la sécurité de l'infrastructure
informatique (performance globale, intégrité de l'environnement). Les contrôles
d'application sont centrés sur les processus métier et les transactions spécifiques.

\paragraph{Approche Proactive vs Réactive}
Les contrôles généraux sont principalement \textbf{proactifs} : ils visent à prévenir
les incidents et à garantir que l'infrastructure est sécurisée avant que des problèmes
ne surviennent. Les contrôles d'application peuvent être à la fois \textbf{proactifs et
réactifs} : ils préviennent les erreurs dans les transactions, mais détectent également
les anomalies après qu'elles se sont produites.

\begin{boxresume}
\textbf{Synthèse du chapitre :}
\begin{itemize}
  \item L'audit porte sur un processus (optimisation/valeur) ; le contrôle porte sur
    une opération (exactitude/conformité).
  \item Le risque est une probabilité dans $]0,1[$ qui nécessite une ressource pour
    exister ; quatre méthodes le traitent.
  \item La formule $\mathrm{RMA} = RI \times RNC \times RND$ encadre le risque de
    mission d'audit.
  \item L'audit interne repose sur trois piliers : spécification, séparation et contrôle
    réciproque des tâches.
  \item La chaîne conceptuelle, point de contrôle \textrightarrow{} objectif de
    contrôle \textrightarrow{} critère \textrightarrow{} question (booléenne)
    \textrightarrow{} constatation (positive ou négative), structure toute mission.
  \item Le contradictoire est un droit de la défense ; l'aveu de carence, son exception.
  \item Les contrôles généraux (GITC) sécurisent l'environnement global ; les contrôles
    d'application (IPO) sécurisent les traitements métier. Les deux sont complémentaires
    et divergent en portée, focalisation et posture proactive/réactive.
\end{itemize}
\end{boxresume}
